\chapter{最後散射與重子生成}

\section{向外看，就是向過去看}

在不同距離上，我們看到的是宇宙在不同時期的樣貌。在近處，回溯時間並不太重要。一年前或一百年前的宇宙，和現在的宇宙並沒有太大差異。但到了足夠遠的地方，我們看到的就是較年輕的宇宙。宇宙越年輕，密度越高；密度越高，溫度也越高。

在某個足夠遠的距離上，宇宙的溫度大約和太陽表面相當。它在我們看來並非如此，因為光子在來到這裡的途中已經發生紅移。紅移因子約為一千，表示觀測到的溫度也降低了相同的倍數：
\[
T_{\mathrm{obs}} \sim \frac{T_{\mathrm{emit}}}{1+z},
\qquad
1+z \sim 10^3.
\]\footnote{這是對00:02:20--00:02:39口頭紅移論證的簡要編輯轉錄。}

那就是宇宙變得不透明的地方。宇宙的主要成分是氫，另有少量氦。溫度高於幾千度時，氫會電離。因此，在那個表面之外，存在的不是中性原子，而是離子：主要是自由運動的質子和電子。那裡的光子能量也更高，基本上是 X 射線。如果偶然形成原子，它很快就會再次被碰撞或光子電離。

電離氣體是導電體，而導電體無法讓輻射順利穿透。電荷會對電場作出反應，來回振盪，並重新發射輻射。因此，在那個距離之外，宇宙對我們而言是不透明的。最後散射面就像太陽的可見表面：它不是一道物質構成的牆，而是介質從不透明突然轉為透明的地方。

那裡的粒子絕大多數是光子，數量約為其他粒子的\(10^8\)倍。今天也是如此，儘管現在光子的能量很低，因此在當前的能量密度中並不重要。一旦溫度降到氫的電離溫度以下，宇宙就變得透明，我們也就能看穿它。

我們的直接觀測到此為止。接下來的問題是，已知的原子、核子與粒子物理，還能讓我們追溯到多早以前。

\section{追溯到更早以前}

就在最後散射面內側，除了光子之外，主要的粒子是質子和電子，還有一些氦核和少量氘核。幾乎沒有反質子，也幾乎沒有正電子。但再往更早以前，溫度終究會高到足以讓兩個光子碰撞，產生一個電子和一個正電子：
\[
\gamma + \gamma \leftrightarrow e^- + e^+.
\]

基本的閾值條件是，光子至少必須提供這一對粒子的靜止能量：
\[
E \ge 2m_e c^2.
\]
低於這個閾值，普通的光無法做到這一點。高於閾值，就能產生電子--正電子對，逆反應也可能發生。在這種溫度下，一鍋熱物質會包含各種粒子反應，就像普通物質中會發生化學反應一樣。

溫度剛高於閾值時，光子仍多於電子和正電子。電子的靜止質量需要很高的能量才能產生。但溫度遠高於閾值時，電子的靜止能量相較於動能便不再重要，粒子密度也會變得相近。此時，電子、正電子和光子的密度大致相同，彼此猛烈碰撞。

\begin{figure}[ht]
\centering
\includegraphics[width=0.62\textwidth]{lecture_06_figure_02.png}
\caption{比較電子與質子質量尺度的對產生閾值示意圖。}
\lectureframe{00:18:47}
\end{figure}

接下來的問題是質子會如何變化。質子已經存在，而且當時的密度很高，但熱能還不足以大量產生質子--反質子對。原因很簡單：
\[
m_p \approx 2000\,m_e,
\qquad
E \ge 2m_p c^2
\]
這是高得多的閾值。\editorialnote{圖中較低處的質量標記不如上方的\(2m_e c^2\)清晰，但口頭比較明確指的是電子與質子的閾值。}

這已經比恆星內部熱得多；在恆星內部，反粒子的豐度並不顯著。

但即使如此，也只是一種暫時的描述。溫度高到足以自由產生質子和反質子時，也高到足以把質子撞散成夸克。因此，真正熾熱的粒子湯，更適合描述為由電子、正電子、夸克、反夸克、光子、中微子、膠子及其他粒子種類組成。黑板上的一個簡單示意反應是
\[
e^- + e^+ \to \gamma \to q + \bar q.
\]

\begin{classroomqa}
\lecturetimestamp{00:20:05}
\audiencequestion{是介子，還是前夸克？}
\lecturerresponse{介子由夸克組成。溫度一旦高到足以打碎質子，也會打碎介子。}
\end{classroomqa}

到目前為止，推理的方向是回溯時間。現在把方向倒過來，看看這鍋熾熱而對稱的粒子湯冷卻時會發生什麼。

\section{電中性與對稱的熱宇宙}

在足夠高的溫度下，宇宙中粒子和反粒子的密度都極高。電子與正電子成對出現，夸克與反夸克也成對出現；自然的圖像是，每種粒子和它的反粒子數量幾乎相等。

電中性這件事需要稍加說明。今天的宇宙是電中性的。一個有用的助記方式是：想像一個閉合且有限的宇宙，例如三維球面，而不是我們實際的宇宙。先用二維球面作為模型來思考。電場線只能在電荷處起始或終止。正電荷發出力線，負電荷吸收力線。在閉合球面上，這些力線不能直接逃到無窮遠，而會在對側重新匯聚。因此，作為助記，閉合且有限的宇宙不能帶有淨電荷。我們的宇宙是否真的閉合，並不是重點。重點是，大尺度上的正、負電荷似乎抵消得極為精確。

電中性\emph{並不}表示質子數和反質子數相同。它表示總正電荷等於總負電荷。質子、正電子以及所有其他帶正電的粒子種類，其電荷總和，必須與電子及所有其他帶負電的粒子種類的電荷總和相平衡。

在已知物理學仍可信的範圍內，或許還能再往前推一點，宇宙曾經熱到足以讓粒子和反粒子大量存在。任何能想到的加速器實驗，在早期宇宙的某個階段都有對應的過程。近乎對稱的起始狀態是很自然的猜想：電子與正電子數目相等，夸克與反夸克數目相等；此外，若以質子和反質子代表重子和反重子，則
\[
N_p = N_{\bar p}.
\]

這個猜想很有吸引力，因為物理定律對粒子和反粒子幾乎對稱。正電子的行為和電子非常相似，雖非完全相同，但相差極微。問題就在於這種近乎對稱。

\section{把電影向前播放}

假設宇宙從一個熾熱狀態開始，粒子和反粒子的豐度相等或幾乎相等，接著逐漸冷卻。

只要溫度夠高，質子和反質子就能由高能碰撞產生，也能湮滅成光子。但溫度終究會降到質子產生閾值以下。此時，質子--反質子對仍能湮滅成光子；然而，宇宙膨脹一段時間後，這些光子損失了足夠的能量，就無法再轉變成質子--反質子對。這就成了單行道。之後，電子和正電子也會發生同樣的事；當時的溫度遠低於質子閾值，但仍遠高於恆星內部的溫度。

在這個過程中，宇宙膨脹的速度遠低於微觀反應速率，因此大部分時間都接近熱平衡。這裡所說的絕熱膨脹，指的就是這種情況。更精確地說，熵的變化非常小。對於熾熱的粒子湯，熵大致等於粒子數，兩者相差的因子約為一。因此，湮滅掉的質子、反質子、電子和正電子，大致由光子取代。

如果質子和反質子的數目完全相等，它們就會全部湮滅。電子和正電子也會如此，最後只剩下光子。但如果質子數哪怕只比反質子數多一點，成對的粒子會湮滅，未配對的多餘部分則會留下來。舉例來說，如果有一百個質子和五十個反質子，五十個反質子會各自找到一個質子，與之湮滅，最後留下五十個質子。電中性也要求電子相對於正電子留下同樣形式的多餘數量。

由今天的光子數，可以推知早期的不對稱程度。今天的光子數約為質子數或電子數的\(10^9\)到\(10^{10}\)倍：
\[
\frac{N_\gamma}{N_p} \sim 10^9\text{--}10^{10}.
\]
所以，在早期熾熱的粒子湯中，大約每一對質子--反質子就多出一個質子。原子之所以至今仍然存在，全靠這微小的不平衡。

\begin{classroomqa}
\lecturetimestamp{00:42:25}
\audiencequestion{為什麼不從統計角度來看這個多出的部分，把它視為冷卻過程中以極小機率產生的結果？}
\lecturerresponse{若要從初始對稱的狀態中多出一個質子，就必須在不產生反質子的情況下產生質子；同時還必須產生與它相伴的負電荷，例如電子。已知的標準模型並沒有這種低能過程。在那些經過檢驗的能量範圍內，重子數守恆，因此一般的冷卻過程不會產生淨多出的重子。}
\end{classroomqa}

所以電荷並不是障礙。示意例子如下：
\[
\gamma + \gamma \to p + e^-.
\]\footnote{這是對00:44:16--00:46:35口頭例子的編輯轉錄，目的是單獨說明：電荷可以守恆，而重子數可以不守恆。}
整體而言，這會是電中性的。但它違反了另一條經驗規則。在已檢驗的標準模型中，產生質子時不能不產生反質子。就目前的討論而言，相關的守恆量是重子數：
\[
B=(N_p+N_n)-(N_{\bar p}+N_{\bar n}).
\]\footnote{這個簡要公式將00:45:18--00:45:40處的口頭定義形式化。}
在已知能量範圍內，重子數守恆，因此這種不對稱不可能在冷卻過程後期、也就是物理學已充分經過檢驗的階段中產生。

\section{重子生成}

這個問題稱為重子生成。假設宇宙起初的重子和反重子數目確實相等。物理定律必須允許什麼，才能讓這個對稱狀態偏向一方，形成觀測到的微小多餘量？

黑板上將問題歸結為兩個符號表示式：
\[
N_p=N_{\bar p},
\qquad
N_p-N_{\bar p},
\]
接著列出一系列條件，說明如何從第一個狀態轉變為第二個狀態。

\begin{figure}[ht]
\centering
\includegraphics[width=0.72\textwidth]{lecture_06_figure_04.png}
\caption{由初始對稱狀態產生重子不對稱的條件。}
\lectureframe{01:22:26}
\end{figure}

\editorialnote{這個畫面中的編號列表只有部分文字清晰可辨。以下三項是根據隨附的口頭說明整理而成。}

第一個條件是重子數不守恆。如果重子數是精確的守恆量，那麼起初淨重子數為零的宇宙，就會一直維持淨重子數為零。

第二個條件是粒子和反粒子之間必須有所偏向。只有重子數不守恆還不夠。如果自然界允許產生質子的過程，以及與之對應、產生反質子的過程，兩者發生的機率相同，平均的不平衡仍然是零。兩者的反應速率必須不同。

一般的統計漲落也遠遠不夠。如果粒子總數約為\(10^{90}\)，其中包括光子，那麼大小約為\(\sqrt{N}\)的漲落只能產生大約\(10^{45}\)，遠不足以解釋宇宙中約有\(10^{80}\)個質子。

第三個條件是，宇宙必須曾在某個時期充分偏離平衡。這裡的論證用到 CPT 定理。如果一團熱氣體永遠處於平衡，倒著播放它的演化，和正著播放看起來會一樣。在這種情況下，結合電荷共軛、宇稱和時間反演的精確對稱性，會有效地恢復粒子與反粒子之間的對稱。要讓真正的偏向留下來並固定下來，實際的歷史在微觀時間尺度上，正向與反向演進時必須有所不同。在宇宙早期，膨脹速度曾足夠快，因此能做到這一點。

因此，三個條件是：
\begin{enumerate}
\item 重子數不守恆；
\item 粒子與反粒子之間的對稱性遭到破壞，更精確地說，是 CP 對稱性破壞；
\item 曾有一個偏離熱平衡的階段。
\end{enumerate}

\begin{classroomqa}
\lecturetimestamp{01:09:08}
\audiencequestion{對照 CPT 定理來看，條件2確實指 CP 破壞，而條件3確實指時間方向的問題嗎？}
\lecturerresponse{是的。條件2確實應該表述為 CP 破壞。條件3指的是實際歷史曾偏離平衡；它不是一條獨立的守恆定律，也不等同於重子數不守恆。}
\end{classroomqa}

有了這三個條件，即使宇宙起初的質子和反質子數目相等，也可能產生偏向其中一方的不平衡。尚未確定的是，自然界是否真的存在重子數不守恆；如果存在，具體機制又是什麼。

\section{已知與未知}

尚未解決的部分是重子數不守恆。我們從未在自然界中直接觀測到會改變
\[
(N_p+N_n)-(N_{\bar p}+N_{\bar n}).
\]
的反應。這可能只是表示，相關反應所需的溫度遠高於實驗室直接測試過的任何溫度；也可能是因為它們在低能下發生的機率極低。但目前這仍是個問號。

這就是質子衰變在這場討論中如此重要的原因。如果重子數不是精確守恆量，原則上質子就應該不穩定。但質子極其穩定。此處引用的壽命下限約為
\[
\tau_p \gtrsim 10^{34}\ \text{years}.
\]
這個結果來自大型實驗：實驗使用大量裝設光偵測器的水，並將設備深埋在地下。它不能證明重子數不會不守恆，只表示任何此類效應在低能下都必須受到極強的抑制。

討論中也提到一項已知物理。在\(100\)到\(250\,\mathrm{GeV}\)附近，弱相互作用會變得更像電磁相互作用。這對詳細模型可能很重要，但並不是當前論證中的關鍵轉折。

暗物質不會大幅改變基本的論述。在多數理論中，暗物質與普通物質的作用非常微弱；就這個問題而言，它基本上只是旁觀者。

\begin{classroomqa}
\lecturetimestamp{01:24:49}
\audiencequestion{引力如何影響這一切？}
\lecturerresponse{主要有兩方面。第一，若沒有宇宙膨脹，就不會出現偏離平衡的階段。第二，黑洞蒸發有力地暗示，重子數不可能永遠是精確的守恆量。除此之外，粒子間一般的引力作用太弱，在這裡影響不大。}
\end{classroomqa}

接著有一個關於電荷的問題：同樣的黑洞論證是否也會威脅電荷守恆？答案是否定的。帶電黑洞具有向外延伸的電通量線，而這些線保留著電荷的資訊。這種黑洞不會完全蒸發，而會逐漸接近極端黑洞的殘餘態。因此，黑洞論證威脅的是重子數守恆，不是電荷守恆。

另一個問題是，極度緻密的早期夸克湯是否應該直接視為黑洞。答案是否定的。關鍵不在於一團靜止的緻密物質；當時的宇宙正在膨脹。黑洞的例子是用來思考重子數可能的命運，並不是說早期宇宙實際上是個受困區域。

\begin{classroomqa}
\lecturetimestamp{01:31:14}
\audiencequestion{人們知道如何修改標準模型，納入這類反應嗎？}
\lecturerresponse{問題恰好相反。要加入過多的重子數不守恆效應很容易，結果會使質子衰變得太快。真正的難題是，既要讓早期宇宙中所需的過程得以發生，又要讓質子維持目前這樣的穩定性。}
\end{classroomqa}

\section{熵與後期的中子問題}

現在回頭整理一下。在絕熱膨脹的宇宙中，熵近似守恆。對於由基本粒子組成的熱粒子湯，熵大致等於粒子數，兩者相差的因子約為一。因此，這個討論不需要假設宇宙後期憑空產生大量光子。湮滅主要是把粒子湯的成分重新分配成輻射。

這個問題的舊版本把因果關係說反了。當時問的不是為什麼相較於光子只有這麼少的質子存留下來，而是反過來問：為什麼相較於質子有這麼多光子？但更精確的問題不是宇宙為什麼在後期產生這麼多熵，而是粒子與反粒子的初始不對稱為什麼這麼小：大約每一對質子--反質子中，只多出一個質子。

自由中子確實會衰變，但這個過程\emph{並不}違反重子數守恆：
\[
n \to p + e^- + \bar\nu.
\]
中子的重子數和質子相同。

自由中子衰變時，可用的衰變能量很小，因為中子的質量只比末態的質量稍大。這也是它的衰變速度以基本粒子的標準來看相當緩慢的原因之一。中子若束縛在原子核中，結合能可能完全抵銷這一點可用能量，於是束縛中子便能保持穩定。

\begin{classroomqa}
\lecturetimestamp{01:43:03}
\audiencequestion{中子只是從宇宙中消失，還是也會有中子產生？}
\lecturerresponse{首先，自由中子會因衰變而減少。但這不是全部情況。中子確實也會透過某些過程產生，例如恆星中的核反應和中子星的形成。因此，中子的總數並非只會一路減少。}
\end{classroomqa}